%% Les trois symboles d'appui d'une poutre : appui simple (sur rouleaux),
%% articulation (appui simple fixe) et encastrement.
\documentclass[tikz,border=4pt]{standalone}
\usepackage{liaisons}
\begin{document}
\begin{tikzpicture}[x=1cm,y=1cm,
  poutre/.style={line width=2.6pt, solideB, line cap=round},
  triangle/.style={line width=1pt, draw=solideE, fill=solideE!12, line join=round},
  titre/.style={font=\small\bfseries, anchor=south},
  legende/.style={font=\small, anchor=north, align=center}]

%% ---------------------- 1. appui simple -----------------------------
\begin{scope}[shift={(0,0)}]
  \draw[poutre] (-0.9,0) -- (0.9,0);
  \draw[triangle] (0,-0.03) -- (-0.42,-0.62) -- (0.42,-0.62) -- cycle;
  \draw[line width=0.9pt, draw=solideE, fill=white] (-0.22,-0.735) circle (0.115);
  \draw[line width=0.9pt, draw=solideE, fill=white] ( 0.22,-0.735) circle (0.115);
  \hachures{(-0.9,-1.07) rectangle (0.9,-0.85)}
  \node[titre] at (0,0.72) {Appui simple};
  \node[legende] at (0,-1.2) {rotation\\ \textbf{et} translation};
\end{scope}

%% ---------------------- 2. articulation -----------------------------
\begin{scope}[shift={(2.9,0)}]
  \draw[poutre] (-0.9,0) -- (0.9,0);
  \draw[triangle] (0,-0.03) -- (-0.42,-0.62) -- (0.42,-0.62) -- cycle;
  \hachures{(-0.9,-0.84) rectangle (0.9,-0.62)}
  \node[titre] at (0,0.72) {Articulation};
  \node[legende] at (0,-1.2) {rotation\\ seulement};
\end{scope}

%% ---------------------- 3. encastrement -----------------------------
\begin{scope}[shift={(5.8,0)}]
  \draw[poutre] (-0.52,0) -- (0.9,0);
  \hachures{(-0.78,-0.62) rectangle (-0.52,0.62)}
  \node[titre] at (0.05,0.72) {Encastrement};
  \node[legende] at (0.05,-1.2) {aucun degré\\ de liberté};
\end{scope}
\end{tikzpicture}
\end{document}
